\documentclass[11pt,a4paper]{article}

%% PACHETE MATEMATICE
\usepackage{amsmath,amssymb,amsfonts}
\usepackage{amsthm}
\usepackage{mathpazo}
\usepackage{eulervm}
\usepackage{enumerate}
\usepackage{color}
\usepackage{graphicx}
\usepackage[all]{xy}
\usepackage{xcolor}
\usepackage{framed}
\usepackage[unicode]{hyperref}
\setcounter{MaxMatrixCols}{20}
\usepackage{multicol}
% \usepackage[romanian]{babel}


%% ASPECT
\linespread{1.05}
\usepackage[T1]{fontenc}
\usepackage[utf8x]{inputenc}
\usepackage[margin=2cm]{geometry}
% \usepackage[color]{showkeys}
\usepackage{anyfontsize}
\usepackage{titlesec}
\titleformat*{\section}{\large\bfseries}

\usepackage{fancyhdr}
\pagestyle{fancy}
\rhead{\small \color{gray} Notițe de Adrian Manea}
\lhead{\small \color{gray} IntProb-FIA, 2017---2018, L. Lemnete \& M. Olteanu}


\usepackage[autostyle = false, style = german]{csquotes}
\usepackage[romanian]{babel}
\newcommand\qq{\enquote}

%% COMENZILE MELE
\renewcommand*{\proofname}{Dem.:}
\newcommand\bb{\mathbb}
\newcommand\dr{\mathrm}
\newcommand\kal{\mathcal}
\newcommand\fk{\mathfrak}
\newcommand\op{\oplus}
\newcommand\ot{\otimes}
\newcommand\ds{\displaystyle}
\newcommand\id{\indent}
\newcommand\nid{\noindent}
\newcommand\seq{\subseteq}
\newcommand\sneq{\subsetneq}
\newcommand\speq{\supseteq}
\newcommand\spneq{\supsetneq}
\newcommand\rar{\rightarrow}
\newcommand\Rar{\Rightarrow}
\newcommand\xrar{\xrightarrow}
\newcommand\lar{\leftarrow}
\newcommand\Lar{\Leftarrow}
\newcommand\xlar{\xleftarrow}
\newcommand\lrar{\leftrightarrow}
\newcommand\Lrar{\Leftrightarrow}
\newcommand\ideal{\trianglelefteq}
\newcommand\ai{\text{ a.\^{i}. }}
\newcommand\sk{(\textit{Schi\c{t}\u{a}}) }
\newcommand\rhup{\rightharpoonup}
\newcommand\rhdn{\rightharpoondown}
\newcommand\lhup{\leftharpoonup}
\newcommand\lhdn{\leftharpoondown}
\newcommand\vid{\emptyset}
\newcommand\wh{\widehat}
\newcommand\ol{\overline}
\newcommand\NN{\mathbb{N}}
\newcommand\RR{\mathbb{R}}
\newcommand\ZZ{\mathbb{Z}}
\newcommand\QQ{\mathbb{Q}}
\newcommand\CC{\mathbb{C}}

%% STILURILE MELE
\newtheoremstyle{dotless}{}{}{\itshape}{}{\bfseries}{:}{ }{}

\theoremstyle{dotless}
\newtheorem{theorem}{Teorem\u{a}}[section]
\newtheorem{proposition}{Propozi\c{t}ie}[section]
\newtheorem{lemma}{Lem\u{a}}[section]
\newtheorem{corollary}{Corolar}[section]
\newtheorem{exercise}{Exerci\c{t}iu}

\newtheoremstyle{dotlessdr}{}{}{}{}{\bfseries}{:}{ }{}
\theoremstyle{dotlessdr}
\newtheorem{example}{Exemplu}[section]
\newtheorem{definition}{Defini\c{t}ie}[section]
\newtheorem{remark}{Observa\c{t}ie}[section]




\begin{document}

\begin{center}
  {\large\textbf{Seminar 9 \\
      Integrale de suprafață}}
\end{center}

\vspace{2cm}

\section*{Noțiuni teoretice}

\indent\indent Integralele de suprafață sînt similare celor curbilinii, cu modificarea
că acum suportul pe care integrăm este o suprafață bidimensională.

Analogul curbelor parametrizate $ \gamma(t) $ din cazul curbiniliu vom avea
\emph{pînze parametrizate} $ \Phi(u, v) $.

\textbf{Integrala de suprafață de speța întîi} se calculează în următorul context
și cu formula de mai jos. Fie $ \Phi : D \seq \RR^2 \to \RR^3 $ o pînză parametrizată,
iar $ \Sigma = \Phi(D) $ imaginea ei. Fie $ F : U \to \RR $ o funcție continuă
pe imaginea pînzei. Integrala de suprafață a lui $ F $ pe $ \Sigma $ este, prin definiție:
\[
  \int_\Sigma F(x, y, z) d \sigma = \iint_D F(\Phi(u,v)) \cdot %
  \Big|\Big| \frac{\partial \Phi}{\partial u} \times \frac{\partial \Phi}{\partial v} \Big|\Big| %
  dudv,
\]
unde $ \Phi(u, v) $ este o parametrizare a suprafeței, iar:
\begin{itemize}
\item $ \vec{T}_u = \frac{\partial \Phi}{\partial u} $ este \emph{vectorul tangent} la
  suprafață în direcția parametrului $ u $;
\item $ \vec{T}_v = \frac{\partial \Phi}{\partial v} $ este \emph{vectorul tangent} la
  suprafață în direcția parametrului $ v $;
\item $ \vec{T}_u \times \vec{T}_v = \vec{N} $ este \emph{vectorul normal} la suprafață.
\end{itemize}

În multe situații, parametrizarea suprafeței $ \Sigma $ este dată într-o formă particulară,
numită \emph{parametrizarea carteziană}. În acest caz, avem $ z = f(x, y) $, cu
$ f : D \seq \RR^2 \to \RR $. Pentru această situație, formula de mai sus se simplifică
și este:
\[
  \int_\Sigma F(x, y, z) d\sigma = \iint_D F(x, y, f(x, y)) \sqrt{1 + p^2 + q^2} dxdy,
\]
unde coeficienții sînt $ p = z_x, q = z_y $.

\vspace{1cm}

Cazuri particulare de interes:
\begin{itemize}
\item Pentru $ F = 1 $, obținem \emph{aria suprafeței} $ \Sigma $;
\item Dacă $ F \geq 0 $ este densitatea unei plăci $ \Sigma $, atunci \emph{masa ei} se
  calculează prin $ M = \ds\int_\Sigma F d\sigma $, iar \emph{coordonatele centrului de %
    greutate} sînt $ x_i^G = \frac{1}{M} \ds\int_\Sigma x_i F d\sigma $;
\item Dacă $ F $ este interpretat ca un cîmp vectorial, atunci integrala de suprafață
  dă \emph{fluxul cîmpului}:
  \[
    \Phi_\Sigma(\vec{F}) = \iint_D F(x(u, v), y(u, v), z(u, v)) \cdot %
    (\vec{T}_u \times \vec{T}_v) dudv.
  \]
\end{itemize}

\vspace{1cm}

Pentru \textbf{integralele de suprafață de speța a doua}, considerăm o
\emph{2-formă diferențială}:
\[
  \omega = Pdy \wedge dz + Q dz \wedge dx + R dx \wedge dy
\]
și fie o pînză parametrizată:
\[
  \Phi : D \seq \RR^2 \to \RR^3, \quad \Phi(u, v) = (X(u, v), Y(u, v), Z(u, v)).
\]

Atunci integrala pe suprafața orientată $ \Sigma $ a formei diferențiale $ \omega $
este definită prin:
\[
  \int_\Sigma \omega = \iint_D \Big( (P \circ \Phi) \cdot \frac{D(Y,Z)}{D(u,v)} + %
  (Q \circ \Phi) \cdot \frac{D(Z, X)}{D(u,v)} + %
  (R \circ \Phi) \cdot \frac{D(X, Y)}{D(u, v)} \Big) dudv,
\]
unde $ \frac{D(Y, Z)}{D(u, v)} $ etc.\ sînt jacobienii funcțiilor $ X, Y, Z $ în
raport cu variabilele $ u, v $. Mai precis, avem, de exemplu:
\[
  \frac{D(Y, Z)}{D(u, v)} =
  \begin{vmatrix}
    Y_u & Y_v \\
    Z_u & Z_v
  \end{vmatrix}.
\]

Putem scrie formula de mai sus într-o variantă simplificată, anume:
\[
  \int_\Sigma \omega = \iint_D %
  \begin{vmatrix}
    P\circ \Phi & Q \circ \Phi & R \circ \Phi \\
    X_u & Y_u & Z_u \\
    X_v & Y_v & Z_v
  \end{vmatrix}
  dudv.
\]


%%%%%%%%%%%%%%%%%%%%%%%%%%%%%%%%%%%%%%%%%%%%%%%%%%%%%%%%%%%%%%%%%%%%%%

\section*{Exerciții}

1. În fiecare din exemplele următoare, considerăm:
\[
  D \ni (u,v) \mapsto \Phi(u,v) = (X(u,v), Y(u,v), Z(u,v)) \in \RR^3
\]
o pînză parametrizată. Să se calculeze vectorii tangenți la suprafață și
versorul normalei la suprafață.

\begin{enumerate}[(a)]
\item \emph{Sfera}: Fie $R > 0$ și $\Phi : [0, \pi] \times [0, 2\pi) \to \RR^3$:
  \[
    \Phi(\theta, \varphi) = (R\sin \theta \cos \varphi, R \sin \theta \sin \varphi, R \cos \theta);
  \]
\item \emph{Paraboloidul}: Fie $a > 0, h > 0$ și $\Phi : [0, h] \times [0, 2\pi) \to \RR^3$:
  \[
    \Phi(u, v) = (au\cos v, au\sin v, u^2);
  \]
\item \emph{Elipsoidul}: Fie $a, b, c > 0$ și $\Phi : [0, \pi] \times [0, 2\pi) \to \RR^3$:
  \[
    \Phi(\theta, \varphi) = (a \sin\theta \cos \varphi, b \sin \theta \sin \varphi, c \cos \theta);
  \]
\item \emph{Conul}: Fie $h > 0, \Phi : [0, 2\pi) \times [0, h] \to \RR^3$:
  \[
    \Phi(u, v) = (v \cos u, v \sin u, v);
  \]
\item \emph{Cilindrul}: Fie $a > 0, 0 \leq h_1 \leq h_2, \Phi : [0, 2\pi) \times [h_1, h_2] \to \RR^3$:
  \[
    \Phi(\varphi, z) = (a \cos \varphi, a \sin \varphi, z);
  \]
\item \emph{Torul}: Fie $0 < a < b, \Phi [0, 2\pi) \times [0, 2\pi) \to \RR^3$:
  \[
    \Phi(u, v) = ((a + b\cos u) \cos v, (a + b \cos u) \sin v, b \sin u).
  \]
\item \emph{Parametrizarea carteziană}: Fie $D \seq \RR^2$ și $f : D \to \RR$, cu
  $f \in \kal{C}^1(D)$:
  \[
    \Phi : D \to \RR^3, \quad \Phi(x, y) = (x, y, f(x,y)).
  \]
\end{enumerate}

\emph{Indicație:} Vectorii tangenți la suprafață sînt $\frac{\partial \Phi}{\partial u}$ și
$\frac{\partial \Phi}{\partial v}$, iar versorul normalei este:
\[
  \frac{1}{\Big|\Big| \frac{\partial \Phi}{\partial u} \times \frac{\partial \Phi}{\partial v}\Big|\Big|}
    \cdot \frac{\partial \Phi}{\partial u} \times \frac{\partial \Phi}{\partial v}.
\]

\vspace{1cm}

2. Fie $\omega = Pdy \wedge dz + Qdz \wedge dx + Rdx \wedge dy$ o 2-formă diferențială,
iar $\Sigma$, imaginea unei pînze parametrizate. Să se calculeze integrala de
suprafață $\int_\Sigma \omega$ în cazurile:
\begin{enumerate}[(a)]
\item $\omega = ydy \wedge dz + zdz \wedge dx + xdx \wedge dy$, unde:
  \[
    \Sigma : \begin{cases}
      X(u, v) &= u \cos v \\
      Y(u, v) &= u \sin v \\
      Z(u, v) &= cv
    \end{cases}
  \]
  cu domeniul $(u, v) \in [a, b] \times [0, 2\pi)$;
\item $\omega = xdy \wedge dz + ydz \wedge dx + zdx \wedge dy$, unde:
  \[
    \Sigma : x^2 + y^2 + z^2 = R^2;
  \]
\item $\omega = yzdy \wedge dz + zxdz \wedge dx + xydx \wedge dy$, unde:
  \[
    \Sigma : \frac{x^2}{a^2} + \frac{y^2}{b^2} + \frac{z^2}{c^2} = 1;
  \]
\item $\omega = xdy \wedge dz + ydz \wedge dx$, unde:
  \[
    \Sigma: x^2 + y^2 = z^2, \quad z \in [1,2];
  \]
\item $\omega = (y + z) dy \wedge dz + (x + y) dx \wedge dy$, unde:
  \[
    \Sigma: x^2 + y^2 = a^2, \quad a \in [0, 1]
  \]
\end{enumerate}

\vspace{1cm}

3. Să se calculeze integrala de suprafață de prima speță:
\[
  \iint_\Sigma F(x, y, z) d\sigma
\]
în următoarele cazuri:
\begin{enumerate}[(a)]
\item $F(x, y, z) = |xyz|$, iar $\Sigma: x^2 + y^2 = z^2, z \in [0, 1]$;
\item $F(x, y, z) = y\sqrt{z}$, iar $\Sigma: x^2 + y^2 = 6z, z \in [0, 2]$;
\item $F(x, y, z) = z^2$, iar $\Sigma = \{(x, y, z) \mid z = \sqrt{x^2 + y^2}, %
  x^2 + y^2 - 6y \leq 0 \}$.
\end{enumerate}

\vspace{1cm}

4. Folosind integralele de suprafață, să se calculeze ariile suprafețelor:
\begin{enumerate}[(a)]
\item sfera de rază $R$;
\item conul $z^2 = x^2 + y^2, z \in [0, h]$;
\item paraboloidul $z = x^2 + y^2, z \in [0, h]$.
\end{enumerate}

\vspace{1cm}

5. Să se calculeze aria suprafeței $\Sigma$ în următoarele cazuri:
\begin{enumerate}[(a)]
\item $\Sigma : 2z = 4 - x^2 - y^2, z \in [0, 1]$;
\item $\Sigma$ este submulțimea de pe sfera $x^2 + y^2 + z^2 = 1$, situată în interiorul
  conului $x^2 + y^2 = z^2$, în semispațiul $z \geq 0$;
\item $\Sigma$ este submulțimea de pe sfera $x^2 + y^2 + z^2 = R^2$, situată în interiorul
  cilindrului $x^2 + y^2 - Ry = 0$;
\item $\Sigma$ este submulțimea de pe paraboloidul $z = x^2 + y^2$, situată în interiorul
  cilindrului $x^2 + y^2 = 2y$.
\end{enumerate}

\vspace{1cm}

6. Să se calculeze fluxul cîmpului vectorial $\vec{V}$ prin suprafața $\Sigma$
în următoarele cazuri:
\begin{enumerate}[(a)]
\item $\vec{V} = x\vec{i} + y\vec{j} + z \vec{k}$, iar $\Sigma: z^2 = x^2 + y^2$,
  cu $z \in [0, 1]$;
\item $\vec{V} = y\vec{i} + x \vec{j} + z^2 \vec{k}$, iar $\Sigma: z = x^2 + y^2$, cu
  $z \in [0, 1]$;
\item $\vec{V} = \frac{1}{\sqrt{x^2 + y^2}} (y\vec{i} - x \vec{j} + \vec{k})$, iar
  $\Sigma: z = 4 - x^2 - y^2, z \in [0, 1]$.
\end{enumerate}





\end{document}
